\documentclass[11pt]{article}

\usepackage[T1]{fontenc}
\usepackage[utf8]{inputenc}
\usepackage{lmodern}
\usepackage[margin=1.04in]{geometry}
\usepackage{amsmath,amssymb,mathtools,amsthm}
\usepackage{aliascnt}
\usepackage{mathrsfs}
\usepackage{microtype}
\usepackage{enumitem}
\usepackage{xcolor}
\usepackage{booktabs}
\usepackage{url}
\usepackage{hyperref}
\usepackage[capitalize,noabbrev]{cleveref}

\hypersetup{
  colorlinks=true,
  linkcolor=blue!45!black,
  citecolor=green!35!black,
  urlcolor=blue!55!black,
  pdftitle={Relative Exponential Periods on Polynomial Chains: Algebraic Independence over Fixed Boundary Fields},
  pdfauthor={Christopher D. Long; Antoine-Auguste Le Blanc},
  pdfsubject={Relative exponential periods, polynomial chains, Fourier--Laplace transform, actual boundary systems, algebraic independence, algebraic specialization},
  pdfkeywords={exponential period, polynomial chain, Fourier--Laplace transform, differential Galois group, relative exactness, E-function, Kontsevich--Zagier}
}

\numberwithin{equation}{section}
\setlength{\emergencystretch}{3em}

\newtheorem{theorem}{Theorem}[section]

\newaliascnt{proposition}{theorem}
\newtheorem{proposition}[proposition]{Proposition}
\aliascntresetthe{proposition}

\newaliascnt{lemma}{theorem}
\newtheorem{lemma}[lemma]{Lemma}
\aliascntresetthe{lemma}

\newaliascnt{corollary}{theorem}
\newtheorem{corollary}[corollary]{Corollary}
\aliascntresetthe{corollary}

\theoremstyle{definition}
\newaliascnt{definition}{theorem}
\newtheorem{definition}[definition]{Definition}
\aliascntresetthe{definition}

\newaliascnt{example}{theorem}
\newtheorem{example}[example]{Example}
\aliascntresetthe{example}

\theoremstyle{remark}
\newaliascnt{remark}{theorem}
\newtheorem{remark}[remark]{Remark}
\aliascntresetthe{remark}

\newcommand{\C}{\mathbb C}
\newcommand{\Q}{\mathbb Q}
\newcommand{\Qbar}{\overline{\mathbb Q}}
\newcommand{\R}{\mathbb R}
\newcommand{\Z}{\mathbb Z}
\newcommand{\N}{\mathbb N}
\newcommand{\Aone}{\mathbb A^1}
\newcommand{\Atwo}{\mathbb A^2}
\newcommand{\Gm}{\mathbb G_m}
\newcommand{\Spec}{\operatorname{Spec}}
\newcommand{\Span}{\operatorname{span}}
\newcommand{\trdeg}{\operatorname{trdeg}}
\newcommand{\per}{\operatorname{per}}
\newcommand{\ord}{\operatorname{ord}}
\newcommand{\End}{\operatorname{End}}
\newcommand{\Hom}{\operatorname{Hom}}
\newcommand{\Ext}{\operatorname{Ext}}
\newcommand{\Gal}{\operatorname{Gal}}
\newcommand{\Lie}{\operatorname{Lie}}
\newcommand{\Frac}{\operatorname{Frac}}
\newcommand{\Sym}{\operatorname{Sym}}
\newcommand{\SL}{\operatorname{SL}}
\newcommand{\GL}{\operatorname{GL}}
\newcommand{\Pic}{\operatorname{Pic}}
\newcommand{\Div}{\operatorname{Div}}
\newcommand{\Res}{\operatorname{Res}}
\newcommand{\Tr}{\operatorname{Tr}}
\newcommand{\Crit}{\operatorname{Crit}}
\newcommand{\supp}{\operatorname{supp}}
\newcommand{\rank}{\operatorname{rank}}
\newcommand{\Ecal}{\mathcal E}
\newcommand{\Fcal}{\mathcal F}
\newcommand{\Hcal}{\mathcal H}
\newcommand{\Lcal}{\mathcal L}
\newcommand{\Mcal}{\mathcal M}
\newcommand{\Ocal}{\mathcal O}
\newcommand{\Pcal}{\mathcal P}
\newcommand{\Rcal}{\mathcal R}
\newcommand{\Vcal}{\mathcal V}
\newcommand{\dd}{\,d}
\newcommand{\KZ}{\mathrm{KZ}}

\title{Relative Exponential Periods on Polynomial Chains:\\
Algebraic Independence over Fixed Boundary Fields}
\author{Christopher D. Long \and Antoine-Auguste Le Blanc\thanks{Le Blanc cryptographic author identifier: \texttt{b6048125\allowbreak{}30b3535b\allowbreak{}23d6dbb0\allowbreak{}f5abe32d\allowbreak{}1ed1a9de\allowbreak{}35fb4f7d\allowbreak{}04d62cff\allowbreak{}08ea19c8}}}
\date{September 2026}


\newcommand{\Hq}{H_q^2}
\newcommand{\M}{\mathscr M}
\newcommand{\Ncal}{\mathcal N}
\newcommand{\Tcal}{\mathcal T}
\newcommand{\A}{\mathbb A}

\newcommand{\PGL}{\operatorname{PGL}}
\newcommand{\wt}{\operatorname{wt}}
\newcommand{\Bcal}{\mathcal B}

\begin{document}
\maketitle

\begin{center}
\small
\texttt{galizur@gmail.com}
\end{center}

\begin{abstract}
We construct rational differential systems for compact polynomial-chain
exponential integrals, with polynomial twisted de Rham representatives and
literal Stokes forcing from the fixed faces. For a real algebraic phase
$q=\alpha x^a+\beta y^b+q_<$, where $a,b\ge2$ are coprime and $q_<$ has lower
weighted degree, we prove a single-chain comparison theorem under explicit
Morse, critical-value independence, face-separation, and saddle-or-continuation
hypotheses. The $(a-1)(b-1)$ monomial moments are algebraically independent
over the fixed-face function field, and their values are algebraically
independent over the corresponding face-value field at every nonzero
algebraic parameter. Every polynomial amplitude has a unique Laurent Stokes
remainder; algebraicity of its value over the face-value field is equivalent
to polynomial twisted exactness.

The proof combines global Picard--Lefschetz geometry, a Fourier
order-versus-degree irreducibility argument, direct simple-spectrum formal
diagonalization, and an elementary stable-lattice obstruction to boundary
adjoint factors. An actual one-column criterion converts these facts into
algebraic independence. We verify both the original elliptic triangle and a
rank-six triangle with phase $x^4+y^3+xy-y$. The boundary fields contain only
the specified faces. No genericity assertion, several-chain comparison,
or general formal-completeness theorem is claimed.
\end{abstract}

\medskip
\noindent\textit{2020 Mathematics Subject Classification.}
Primary 11J91; Secondary 11J81, 12H05, 32S40, 34M35.
\tableofcontents

\section{Introduction and main comparison theorem}\label{sec:intro}
We study actual compact integrals and their polynomial Stokes certificates,
not an unspecified completion of a formal period algebra. The basic
polynomial boundary construction is general; the comparison theorem below
verifies its actual-column criterion for a weighted-plane class under
explicit hypotheses. The elliptic and rank-six examples give two complete
applications.

The algebraic changes of variables, linearity, and twisted Stokes identities
used here belong to the exponential Kontsevich--Zagier framework
\cite{KontsevichZagier,CommelinHabeggerHuber}. The conclusions below concern
one actual column over its fixed boundary field. They do not assert a
presentation of all formal exponential-period relations. The functional and
numerical independence proofs do not use the companion formal-comparison
theorems or Paper IV.

\paragraph{Structure and previous version.}
Section~\ref{sec:boundary-construction} constructs actual systems and their
$E$-function coordinates. Section~\ref{sec:criterion} proves the general
one-column criterion. The later sections verify its hypotheses using the
geometry, Fourier irreducibility, and the fixed-face obstruction, and then
specialize at nonzero algebraic parameters. Section~\ref{sec:elliptic-example}
retains the explicit continuation needed for the original triangle.
The pre-integration version \cite{PaperVPrevious}, including its separate
observable-Lefschetz, Hodge, face-lattice, and multicolumn discussions, is
archived unchanged. Those additional discussions are not inputs to the
proof here. The earlier unrestricted endpoint-complete claim remains
withdrawn; it is not restored by this theorem.

\paragraph{Revision.} The 6 October 2026 revision integrates the
source-checked single-chain theorem and its stable-lattice proof. The
manuscript's author list and original month-year date are retained.



Throughout $k=\Qbar$, $k_\R=k\cap\R$, $F_0=k(z)$, $F=\C(z)$.

\begin{definition}[Polynomial two-chain]\label{def:polynomial-chain}
A polynomial two-chain over $k_\R$ is a finite sum
\[
 D=\sum_{\ell=1}^N a_\ell(\Phi_\ell)_*[0,1]^2,
 \qquad a_\ell\in k,\quad
 \Phi_\ell\in k_\R[u,v]^2,
\]
with the standard orientation on each square.  The restrictions of the
$\Phi_\ell$ to the oriented edges of the squares are the \emph{faces} of
$D$; $|\partial D|$ denotes the union of their images in $\R^2$.
Cancellations between coincident oriented faces are allowed.
\end{definition}

For a polynomial phase $q$ and a polynomial two-chain $D$, put
\[
 I_D(\omega)(z)=\int_D e^{zq}\omega,\qquad d_z=d+z\,dq\wedge,\qquad
 \Hq=\Omega^2_{F_0[x,y]/F_0}/d_z\Omega^1_{F_0[x,y]/F_0}.
\]
$K_{\partial,0}$ is the differential field over $F_0$ generated by the
polynomial-amplitude integrals over the fixed faces of $D$ and the
exponentials at their endpoints; $K_\partial=\C K_{\partial,0}$;
$K_{\partial,\xi}$ is the field generated over $k$ by the values at $\xi$.
For a face $e$ with parametrization
$\phi_e:[a_e,b_e]\to\R^2$ put $p_e=q\circ\phi_e$ and
$d_e=\deg p_e$.

\begin{definition}[The class $(W)$]\label{def:W}
Fix coprime integers $a,b\ge2$ and give $x,y$ the weights $\wt(x)=b$,
$\wt(y)=a$. A phase is of class $(W)$ if
\[
 q=\alpha x^a+\beta y^b+q_<,\qquad \alpha\beta\ne0,\qquad
 q_<\in k_\R[x,y]\ \text{of weighted degree}<ab,
\]
with $\alpha,\beta\in k_\R$. Put $\mu=(a-1)(b-1)$ and
$B=\{x^iy^j:0\le i\le a-2,\ 0\le j\le b-2\}$.
\end{definition}

The elliptic phase $y^2-x^3+3x$ is of class $(W)$ with $(a,b)=(3,2)$, $\mu=2$.

\begin{definition}[Chain multiplicity]\label{def:chain-multiplicity}
Let $D=\sum_\ell a_\ell(\Phi_\ell)_*[0,1]^2$ be a polynomial
two-chain.  For $P\notin|\partial D|$ define
\[
 m_D(P)=\sum_\ell a_\ell\,
 \deg(\Phi_\ell,(0,1)^2,P),
\]
where the last factor is the Brouwer degree.  Thus $m_D$ is a compactly
supported, locally constant $k$-valued function on
$\R^2\setminus|\partial D|$.  We call $m_D(P)$ the
\emph{multiplicity of $D$ at $P$}.  This definition includes
non-injective parametrizations.
\end{definition}

\begin{lemma}[Phase densities]\label{lem:phase-density}
Let $q\in k_\R[x,y]$ have finite critical locus, let $D$ be a polynomial
two-chain over $k_\R$, and let $r\in k[x,y]$.
\begin{enumerate}[label=\textup{(\alph*)},leftmargin=2.2em]
\item There is a compactly supported $L^1$ function
$\rho_{D,r}:\R\to\C$, unique almost everywhere, such that
\[
 I_D(r\,dx\wedge dy)(z)
 =\int_\R e^{zt}\rho_{D,r}(t)\,dt.
\]
At every regular value $t$ it has the coarea representative
\[
 \rho_{D,r}(t)=
 \int_{q^{-1}(t)}
 \frac{m_D(x,y)r(x,y)}{\|\nabla q(x,y)\|}
 \,d\mathcal H^1.
\]
On an interval containing no critical value of $q$, no critical value
of a face phase $p_e$ on its parameter interval, and no vertex value,
this representative is real analytic; its germ therefore has a
holomorphic complexification.
\item For a face moment
\[
 J_{e,h}(z)=\int_{a_e}^{b_e}h(s)e^{zp_e(s)}\,ds,
 \qquad h\in k[s],
\]
and an interval avoiding the critical values and endpoint values of
$p_e$, the inverse-Laplace density is algebraic over $\C(t)$.  More
precisely, after splitting into monotone inverse branches $s_j(t)$ it is
\[
 \rho_{e,h}(t)=
 \sum_j \frac{h(s_j(t))}{|p_e'(s_j(t))|}
 =\sum_j \varepsilon_j\,
   \frac{h(s_j(t))}{p_e'(s_j(t))},
 \qquad \varepsilon_j\in\{\pm1\}.
\]
Endpoint exponentials (and the constant function $1$) are Dirac point
masses at their phase values.
\end{enumerate}
\end{lemma}
\begin{proof}
For each summand of $D$, the degree formula for a map of the square gives
\[
 \int_D r e^{zq}\,dx\,dy
 =\int_{\R^2}m_D(x,y)r(x,y)e^{zq(x,y)}\,dx\,dy.
\]
This also proves directly that non-injective parametrizations are counted
with their signed local degrees.  Since $m_D$ has compact support, the
coarea formula applied to $q$ gives (a), including $L^1$-integrability.
Away from the stated exceptional values, the implicit-function theorem
makes the regular level arcs and their intersections with the polynomial
faces vary real analytically.  Splitting the compact support into the
finitely many resulting analytic pieces and integrating in analytic
coordinates gives the asserted real-analytic representative.

For (b), split $[a_e,b_e]$ into monotonicity intervals for $p_e$.
Ordinary one-variable change of variables gives the displayed formula.
Each inverse branch $s_j(t)$ is algebraic, and the sign of $p_e'$ is
constant on the branch, so the density is algebraic.  A value supported
at an endpoint is the stated Dirac mass.
\end{proof}

Whenever we speak below of the density germ of
$I_D(r\,dx\wedge dy)$, we mean the analytic germ of the representative
in Lemma~\ref{lem:phase-density}(a).

\paragraph{Hypotheses.} Let $q$ be of class $(W)$ and let $D$ be a
polynomial two-chain over $k_\R$ in the sense of
Definition~\ref{def:polynomial-chain}.  This arithmetic restriction is
used in (iii).
\begin{enumerate}[label=\textup{(\Alph*)},leftmargin=2.6em]
\item[(M)] \emph{Morse.} $q$ has $\mu$ distinct critical values
 $c_1,\dots,c_\mu$ (equivalently: all critical points are nondegenerate with
 distinct values).
\item[(Q)] \emph{Independent critical values.} $c_2-c_1,\dots,c_\mu-c_1$ are
 linearly independent over $\Q$.
\item[(F)] \emph{Faces.} For every face $e$ with $d_e\ge2$, the set
 $\{c_i-c_j:i\ne j\}$ is not contained in
 $\{u-v: u,v\ \text{critical values of }p_e\}$.
\item[(S)] \emph{Saddle.} Some real critical point $P_0$ of $q$ with
 $\det\operatorname{Hess}q(P_0)<0$ lies off the image of $\partial D$, the
 multiplicity $w:=m_D(P_0)$ is nonzero, and $c:=q(P_0)$ is neither the
 value of $q$ at a vertex of $D$ nor a critical value of any $p_e$ on its
 parameter interval.
\end{enumerate}
For $\mu=2$, (Q) is automatic.

\begin{theorem}[Single-chain comparison]\label{thm:main}
Let $q$ be of class $(W)$ satisfying \textup{(M)} and \textup{(Q)}, and let $D$
be a polynomial two-chain over $k_\R$ satisfying \textup{(F)} and either \textup{(S)} or \textup{(S$^*$)} below. Put
$M_m=I_D(m\dd x\wedge dy)$ for $m\in B$.
\begin{enumerate}[label=\textup{(\roman*)},leftmargin=2.4em]
\item The classes of $m\,dx\wedge dy$, $m\in B$, form an $F_0$-basis of $\Hq$,
 and a $k[z,z^{-1}]$-basis of the lattice of polynomial forms. Every
 $r\in k[x,y]$ has a unique remainder and a polynomial Stokes certificate
 \[
  r=\sum_{m\in B}a_{r,m}(z)\,m+L_xP_r+L_yQ_r,\qquad a_{r,m}\in k[z^{-1}],\quad
  P_r,Q_r\in z^{-1}k[z^{-1}][x,y],
 \]
 with $L_x=\partial_x+zq_x$, $L_y=\partial_y+zq_y$.
\item The $\mu$ functions $M_m$ are algebraically independent over $K_\partial$.
\item For every $\xi\in k^\times$ the $\mu$ numbers $M_m(\xi)$ are algebraically
 independent over $K_{\partial,\xi}$.
\item For $r\in k[x,y]$ and $\xi\in k^\times$, the value $I_D(r\dd x\wedge dy)(\xi)$
 is algebraic over $K_{\partial,\xi}$ iff it lies in $K_{\partial,\xi}$ iff
 $r\,dx\wedge dy\in(d+\xi\,dq\wedge)\Omega^1_{k[x,y]/k}$.
\end{enumerate}
\end{theorem}

\begin{remark}[The continuation alternative]\label{rem:Sstar}
Hypothesis (S) is used only in Lemma~\ref{lem:saddle}. It may be replaced by:
\begin{enumerate}[leftmargin=2.6em]
\item[(S$^*$)] on some interval $e_0\subset\R$ free of critical values of $q$
 and of the $p_e$, of vertex values and of $0$, the density germ
 $\rho=\rho_{D,1}$ of $I_D(dx\wedge dy)$ from
 Lemma~\ref{lem:phase-density} continues analytically along a path to a punctured
 neighbourhood of some $c_i$, where its variation around $c_i$ is a nonzero
 constant multiple of the vanishing period $\int_{\nu_i(t)}dx\wedge dy/dq$.
\end{enumerate}
Lemma~\ref{lem:elliptic-continuation} verifies $(S^*)$ for the original
elliptic triangle: its saddle $(1,0)$ lies outside the triangle, so (S)
itself fails. Section~\ref{sec:elliptic-example} verifies (F) and derives
its functional and numerical comparison from Theorem~\ref{thm:main}.
\end{remark}

\begin{corollary}[Polygons]\label{cor:polygon}
Let $\min(a,b)\ge3$, let $q$ be of class $(W)$ with \textup{(M)}, \textup{(Q)},
and let $D$ be a real polygonal chain (all faces straight segments) with vertices
in $k_\R^2$. Then \textup{(F)} holds automatically. If moreover the coefficients
of $q$ and the vertices lie in a field $k_0\subset k_\R$ over which
$[k_0(c_i):k_0]\ge\max(a,b)$ for all $i$, then \textup{(S)} holds as soon as some
real saddle of $q$ lies in the interior of $D$ with nonzero multiplicity.
\end{corollary}
\begin{proof}
A straight face has $d_e\le\max(a,b)$, so $p_e$ has at most $\max(a,b)-1<\mu$
critical values and at most $(\mu-1)(\mu-2)<\mu(\mu-1)$ nonzero differences,
whereas under (Q) the $\mu(\mu-1)$ numbers $c_i-c_j$ are pairwise distinct.
Vertex values lie in $k_0$ and critical values of $p_e\in k_0[s]$ have degree at
most $d_e-1<\max(a,b)$ over $k_0$, so neither can equal $c$.
\end{proof}

\begin{example}[A rank-six instance]\label{ex:intro}
For
\[
 q=x^4+y^3+xy-y,\qquad D=\operatorname{conv}\{(0,-1),(1,-1),(\tfrac12,0)\},
\]
the six functions $\int_D x^iy^je^{zq}\dd x\dd y$ $(0\le i\le2,\ 0\le j\le1)$ are
algebraically independent over the field of the three edges, and so are their
values at every nonzero algebraic $\xi$ over the edge-value field. See
\S\ref{sec:example}.
\end{example}

\paragraph{Differential-module convention.}
For a system $Y'=EY$ we use its coefficient module: a column of basis
elements $e$ satisfies $\partial e=Ee$. Solutions are differential
homomorphisms out of this module, so the solution functor is contravariant.
Its horizontal representation is dual to the solution representation.
The natural trace duality identifies the corresponding endomorphism and
trace-free endomorphism representations. All module embeddings below are
under the covariant horizontal-fibre equivalence with these identifications.

\section{Polynomial representatives and literal boundary forcing}\label{sec:boundary-construction}
Put $k=\overline{\Q}$ and $F_0=k(z)$.  Let $D$ be a polynomial
two-chain over $k_\R$ as in Definition~\ref{def:polynomial-chain}, and
let $q\in k[x,y]$. For a polynomial two-form
$\omega$, write
\[
 I_D(\omega)=\int_D e^{zq}\omega.
\]
When coefficients of $\omega$ lie in $F_0$, interpret this as a meromorphic
function of $z$ by linearity. For each oriented face parametrization
$\phi_e:[a_e,b_e]\to\A^2$, put $p_e=q\circ\phi_e$.
The actual boundary field $K_{\partial,0}$ is the differential field over
$F_0$ generated by all functions
\[
 \int_{a_e}^{b_e} h(t)e^{zp_e(t)}\dd t,\qquad h\in k[t],
\]
and the exponential functions at their algebraic endpoints. Only the fixed
faces are used; this does not mean the field of all compact line periods.

Let $d$ differentiate only in $x,y$, and set
\[
 d_z=d+z\,dq\wedge,\qquad
 \Hq=\Omega^2_{F_0[x,y]/F_0}/d_z\Omega^1_{F_0[x,y]/F_0}.
\]

\begin{theorem}[Polynomial-representative boundary system]
\label{thm:boundary-system}
The space $\Hq$ is finite dimensional. There are forms
$\omega_i=r_i(x,y)\,dx\wedge dy$, with $r_i\in k[x,y]$, whose classes form
an $F_0$-basis, a matrix $A\in M_n(F_0)$, and polynomial-in-$x,y$ one-forms
$\beta_i\in\Omega^1_{F_0[x,y]/F_0}$ such that
\begin{equation}\label{eq:reduction-system}
 q\omega_i=\sum_jA_{ij}(z)\omega_j+d_z\beta_i.
\end{equation}
The actual integral column $U_i=I_D(\omega_i)$ satisfies
\begin{equation}\label{eq:actual-system}
 U'=AU+b,\qquad b_i=\int_{\partial D}e^{zq}\beta_i\in K_{\partial,0}.
\end{equation}
Every polynomial two-form $\omega$ has an identity
\begin{equation}\label{eq:amplitude-normal}
 \omega=\sum_i a_i(z)\omega_i+d_z\beta,
\end{equation}
where $a_i\in F_0$ and $\beta$ is polynomial in $x,y$. Consequently
\begin{equation}\label{eq:actual-normal}
 I_D(\omega)=\sum_i a_i(z)U_i+
              \int_{\partial D}e^{zq}\beta.
\end{equation}
For any finite set of amplitudes, these functions belong to a finite closed
rational first-order system whose coordinates are actual $E$-functions.
No spatial denominator is introduced. Parameter denominators may occur.
\end{theorem}

\begin{proof}
The polynomial Gauss--Manin complex for $q:\A^2\to\A^1_t$ is
\[
 \bigl(\Omega^{\bullet+2}_{k[x,y]/k}[\partial_t],
             d-dq\wedge\partial_t\bigr),
\]
and its cohomology is holonomic \cite[Proposition~9.1]{Sabbah}.
The Fourier substitution $\partial_t\mapsto-z$, $t\mapsto\partial_z$
turns its differential into $d+z\,dq\wedge$. Fourier transform is exact
on Weyl modules, and localization to $k(z)$ is exact. The degree-zero
cohomology of the localized transformed complex is therefore precisely
$\Hq$, which is finite dimensional over $k(z)$ by holonomicity.
This uses neither cohomological tameness nor the rank and lattice
conclusions stated later in \cite{Sabbah}. In particular, no identification
with a Milnor number is asserted here.

The monomial two-forms with coefficients in $k$ span the quotient over
$F_0$, so a finite subset is a basis. The connection on the quotient is
\[
 \nabla_z[\omega]=[\partial_z\omega+q\omega],
\]
because
\[
 [\partial_z+q,d+z\,dq\wedge]
 =dq\wedge-dq\wedge=0.
\]
Expressing $q\omega_i$ in the basis proves \eqref{eq:reduction-system};
expressing an arbitrary class proves \eqref{eq:amplitude-normal}.

For a one-form polynomial in the spatial variables, there is the literal
identity
\[
 e^{zq}d_z\beta=d_{x,y}(e^{zq}\beta).
\]
Stokes on the original parametrized compact chain proves
\eqref{eq:actual-system} and \eqref{eq:actual-normal}. Each pullback of
$\beta$ to a face is a finite $F_0$-linear combination of polynomial
one-forms in its parameter. Thus every forcing term is an actual face
integral with coefficients in $F_0$.

To see finiteness on the boundary, if $\deg p=d\ge2$, division by
$\partial_t+zp'(t)$ reduces every polynomial amplitude to degree at most
$d-2$. The leading degree of $(\partial_t+zp')h$, for $h\ne0$, is
$\deg h+d-1$, with nonzero leading coefficient in $F_0$. The remainder is
therefore unique. Integration of the exact part gives only endpoint
exponentials. Multiplication by $p$ and the same division give a rational
first-order system for the reduced moments and the endpoint exponentials.
For $d=1$ the twisted operator is surjective on polynomials; for $d=0$ ordinary polynomial antiderivatives suffice. In both cases only endpoint terms remain.
For finitely many faces, this supplies a finite boundary vector; its
fraction field is $K_{\partial,0}$.

The integral coordinates are entire and have algebraic Taylor coefficients.
After pullback to the square they have the form
$\int_{[0,1]^2}P(u,v)e^{zQ(u,v)}\dd u\dd v$. For the first $m$ Taylor
coefficients, one denominator of the form
\[
 d^{m+1}\operatorname{lcm}(1,\ldots,Cm+C_0)^2
\]
suffices, with fixed integers $d,C,C_0$; all conjugates grow exponentially.
The same argument in one variable handles the faces. Together with the
constructed differential systems, these are the $E$-function conditions.
\end{proof}

\begin{proposition}[Actual affine comparison under a change of representatives]
\label{prop:gauge}
Suppose
$\widetilde\omega=T(z)\omega+d_z\gamma$, with $T\in\GL_n(F_0)$.
Put $g=\int_{\partial D}e^{zq}\gamma$. Then the actual columns satisfy
\[
 \widetilde U=TU+g,
\]
and their systems are related by
\[
 \widetilde A=(T'+TA)T^{-1},\qquad
 \widetilde b=Tb+g'-\widetilde A g.
\]
In particular, this is a comparison of the actual affine solution objects,
not a statement only about their homogeneous composition factors.
\end{proposition}
\begin{proof}
The first equality is Stokes, and differentiating it proves the other two.
If two primitives give the same form identity, their difference is
$d_z$-closed and its integral over $\partial D$ is zero by the same Stokes
identity. Thus primitive choices do not introduce an arbitrary constant.
\end{proof}

After adjoining complex constants, write $F=\C(z)$,
$K_\partial=\C K_{\partial,0}$, and work in a field $\M$ of meromorphic
germs containing the actual functions. Integration modulo the actual
boundary field gives a canonical differential map
\begin{equation}\label{eq:canonical-comparison}
 K_\partial\otimes_{F_0}\Hq\longrightarrow\M/K_\partial,
 \qquad [\omega]\longmapsto I_D(\omega)\bmod K_\partial.
\end{equation}
Here $\M/K_\partial$ is used as a differential vector space, not as a
quotient field or ring. This map need not be injective for arbitrary chains.
The next criterion proves injectivity rather than building it into a
completion definition.

\section{An actual-solution comparison criterion}\label{sec:criterion}
Let $\mathcal B$ be a finite rational boundary differential system, let
$L_\partial/F$ be a Picard--Vessiot field for that system containing its
actual solution coordinates, and let $\Ncal\subset L_\partial$ be their
finite-dimensional $F$-linear differential span. Include $1$ when needed and put $K_\partial=F(\Ncal)\subset L_\partial$. All fields are embedded in a common no-new-constants differential overfield.
Let $\Tcal_\partial$ be the tensor category generated by $\mathcal B$,
allowing subquotients and duals. Its objects are precisely the differential
modules trivialized by $L_\partial$ \cite[Theorem 1.5.2]{Singer}.

\begin{lemma}[Differentially finite elements of a Picard--Vessiot field]
\label{lem:dfinite}
Every finite-dimensional $F$-differential subspace of $L_\partial$ belongs
to $\Tcal_\partial$.
\end{lemma}
\begin{proof}
Let $R$ be the Picard--Vessiot ring. If $f\in L_\partial$ satisfies a scalar
linear equation over $F$, all its Galois translates satisfy that equation.
Their constant span is finite dimensional. After trivializing the torsor,
a rational function with finite-dimensional translation orbit has no pole:
the union of the polar loci of a finite basis is a proper closed subset
invariant under all translations, whereas the group acts transitively on
itself. This remains valid component by component if the group is
 disconnected. Thus $f\in R$; see the locally finite argument in
\cite[proof of Proposition 1.3.33]{Singer}.

Write $R=F[Y_{ij},(\det Y)^{-1}]$ for a boundary fundamental matrix.
A bounded span of monomials in $Y_{ij}$ and $(\det Y)^{-1}$ is stable
under differentiation, since $Y'=A_\partial Y$ and
$((\det Y)^{-1})'=-\operatorname{tr}(A_\partial)(\det Y)^{-1}$.
Such a span is a quotient of finite sums of tensor constructions in the
boundary module and its determinant dual. Every finite-dimensional
$F$-differential subspace of $R$ lies in one of these spans. Closure under
subquotients and the Picard--Vessiot tensor equivalence
\cite[Theorem 1.5.2]{Singer} prove the claim.
\end{proof}

\begin{lemma}[Regular affine cocycles]\label{lem:affine-cocycle}
Let $G$ be an algebraic group over an algebraically closed field $C_0$
of characteristic zero, with reductive identity component, and let $V$
be a finite-dimensional rational representation. Every regular cocycle
$c(gh)=c(g)+g c(h)$ has the form $c(g)=v-gv$ for some $v\in V$.
\end{lemma}
\begin{proof}
Let $g$ act on $V\oplus C_0$ by
$g(w,t)=(gw+t c(g),t)$. Complete reducibility supplies an invariant
lift $(v,1)$ of $1$ in the trivial quotient, giving the assertion.
Complete reducibility here includes disconnected groups: average a
$G^\circ$-equivariant splitting over the finite component group
\cite[Corollary~22.43]{MilneAG}. Equivalently, this is the vanishing of
algebraic $H^1(G,V)$ in \cite[Proposition~15.15]{MilneAG}.
\end{proof}

\begin{theorem}[Actual one-column comparison]
\label{thm:criterion}
Suppose an actual column $U$ satisfies $U'=AU+b$, with $b_i\in\Ncal$,
and $n\ge2$. Assume:
\begin{enumerate}
\item the homogeneous Galois group over $F$ contains $\SL_n$;
\item its adjoint module $\End^0(A)$ does not belong to $\Tcal_\partial$;
\item at least one coordinate of $U$ is not in $\Ncal$.
\end{enumerate}
Then $U_1,\ldots,U_n$ are algebraically independent over
$K_\partial\subset L_\partial$. In the setting of
\cref{thm:boundary-system}, \eqref{eq:canonical-comparison} is injective,
and its image is the span of the actual polynomial-amplitude periods
modulo $K_\partial$.
\end{theorem}

\begin{proof}
Let $E=\Ncal+\sum_i FU_i$. The differential module $E/\Ncal$ is a
quotient of the coefficient module associated with $A$. Use the coefficient module with basis $e_i$ and
$\partial e_i=\sum_j A_{ij}e_j$; the map $e_i\mapsto U_i\bmod\Ncal$
is a differential map. Its horizontal representation is dual to the
solution representation of $Y'=AY$, so it is irreducible and has the same
adjoint tensor module. By the third assumption, the quotient is nonzero and hence is
that entire irreducible module.

If all $U_i$ belonged to $L_\partial$, \cref{lem:dfinite} would put $E$,
then $E/\Ncal$, and then its adjoint module in $\Tcal_\partial$,
contradicting the second assumption. Thus $U\notin L_\partial^n$.

In a joint Picard--Vessiot extension for $A$ and the boundary system, the
kernel of restriction to $L_\partial$ maps to a normal subgroup $H$ of
the homogeneous group $G\subset\GL_n$. A normal subgroup of a group
containing $\SL_n$ either contains $\SL_n$ or has trivial projective
image. Indeed, if $H\cap\SL_n$ is not all of $\SL_n$, it is finite central.
For each $h\in H$, the commutator map from the connected group $\SL_n$
to this finite intersection is constant and equals the identity at $1$.
Thus $h$ centralizes $\SL_n$ and is scalar. This also excludes a hidden
noncentral finite projective image. In the latter, projectively trivial case, $H$ acts
trivially on $\End^0(A)$, so the adjoint module is trivialized by
$L_\partial$, a contradiction. Therefore the homogeneous group after
base change to $L_\partial$, and hence to $K_\partial$, contains $\SL_n$.

Let $Y$ be a fundamental matrix over a joint field for the inhomogeneous
system over $K_\partial$. Its affine Galois group is a subgroup of
$\C^n\rtimes G_{K_\partial}$, acting by
$\sigma(U)=U+Yc_\sigma$. Its translation kernel is an invariant vector
subspace, hence either $0$ or $\C^n$.
If it is $\C^n$, translations make the orbit of $U$ Zariski dense in
$\A^n$. If it is zero, the remaining regular algebraic cocycle of the
homogeneous group is a coboundary by \cref{lem:affine-cocycle}.
Here reductivity of its identity component follows
also directly: a unipotent radical would lie in $\SL_n$ and be normal
under $\SL_n$, which has no nontrivial connected normal unipotent
subgroup. Explicitly, writing $\sigma(Y)=Yg_\sigma$, the cocycle
can be written $c_\sigma=g_\sigma c-c$; hence $U-Yc$ is fixed. Thus
\[
 U=a+Yc,\qquad a\in K_\partial^n,\quad c\in\C^n.
\]
The vector $c$ is nonzero, since otherwise $U\in K_\partial^n$. The
$\SL_n$-orbit of a nonzero vector is $\A^n\setminus\{0\}$, again dense.
In either case a polynomial over $K_\partial$ vanishing at $U$ vanishes
on a dense affine orbit and is zero. Linear independence modulo the
boundary field, and then the last assertion by normal-form reduction,
follow immediately.
\end{proof}

\begin{remark}
The third condition concerns only rational \emph{linear} combinations of a
finite actual boundary vector. It is strictly weaker than the desired
nonmembership in the boundary fraction field. The proof is what upgrades
this linear test to the fraction-field statement. A nonzero abstract
cohomology class alone would not suffice.
\end{remark}

\section{Normal form for the class \texorpdfstring{$(W)$}{(W)}}

\begin{lemma}\label{lem:normal}
Let $q$ be of class $(W)$ (coprimality is not used here).
\begin{enumerate}[label=\textup{(\alph*)},leftmargin=2.2em]
\item Every $r\in k[x,y]$ can be written
 $r=\sum_{m\in B}a_m(z)m+L_xP+L_yQ$ with $a_m\in k[z^{-1}]$ and
 $P,Q\in z^{-1}k[z^{-1}][x,y]$.
\item If $\sum_{m\in B}a_m m=L_xP+L_yQ$ with $a_m\in F_0$ and
 $P,Q\in F_0[x,y]$, then all $a_m=0$.
\item Consequently $B$ is an $F_0$-basis of $\Hq$ and a $k[z,z^{-1}]$-basis of
 $G=\Omega^2[z,z^{-1}]/d_z\Omega^1[z,z^{-1}]$. In this basis
 $\nabla_z[m]=[qm]=\sum_{m'}A_{mm'}[m']$ with
 \[
  A=A_0+A_1z^{-1}+\dots+A_sz^{-s}\in M_\mu(k[z^{-1}]),
 \]
 and $A_0$ is the matrix of multiplication by $q$ on
 $k[x,y]/(q_x,q_y)$ in the basis $B$. Under \textup{(M)} its eigenvalues are
 $c_1,\dots,c_\mu$, each simple.
\end{enumerate}
The same statements hold with $k$ replaced by $\C$, and the reduction and independence arguments also apply after specialization at every $z=\xi\ne0$.
\end{lemma}
\begin{proof}
(a) Let $x^my^n$ be a monomial of $r$, with coefficient $\gamma$, of maximal
weight among those not in $B$. If $m\ge a-1$ put $v=\gamma x^{m-a+1}y^n/(a\alpha z)$.
Then $L_xv=\gamma x^my^n+(\gamma/a\alpha)\,q_{<,x}x^{m-a+1}y^n+v_x$, and the last
two terms have weight strictly below $\wt(x^my^n)$ because
$\wt(q_{<,x})<ab-b$. Replace $r$ by $r-L_xv$. If $m<a-1$ then $n\ge b-1$ and one
uses $L_y$ symmetrically. The maximal weight of a monomial outside $B$, with its
multiplicity among the finitely many monomials of that weight, strictly
decreases, so the process stops.

(b) Among all pairs $(P,Q)$ representing the given combination choose one
minimizing $\varpi=\max(\wt P+ab-b,\ \wt Q+ab-a)$. Let $P^{\rm top}$ and $Q^{\rm top}$ be the components of $P$ and $Q$ of weights
$\varpi-ab+b$ and $\varpi-ab+a$ (one of them may be zero). The weight-$\varpi$
component of $L_xP+L_yQ$ is $z(a\alpha x^{a-1}P^{\rm top}+b\beta y^{b-1}Q^{\rm top})$. If it
is nonzero it is a nonzero element of the monomial ideal $(x^{a-1},y^{b-1})$,
which cannot equal the weight-$\varpi$ component of a combination of monomials
of $B$. Hence it vanishes, and since $x^{a-1},y^{b-1}$ is a regular sequence,
$P^{\rm top}=b\beta y^{b-1}R$, $Q^{\rm top}=-a\alpha x^{a-1}R$. As $L_x$ and $L_y$
commute, $(P-L_y(R/z),\,Q+L_x(R/z))$ represents the same element and has smaller
$\varpi$ unless $P=Q=0$. So $P=Q=0$ and all $a_m=0$.

(c) The first two assertions follow from (a) and (b). In (a), the terms of
order $z^0$ give $r\equiv\sum a_m(\infty)m$ modulo $(q_x,q_y)$, and the argument
of (b) with $z$ removed shows $B$ is a basis of the Jacobian algebra; this
identifies $A_0$. That algebra is the product of its local algebras at the
critical points, on each of which $q$ has the single eigenvalue $q(P)$, with
multiplicity the local Milnor number.
\end{proof}

\section{Period calculus and irreducibility of \texorpdfstring{$\Hq$}{Hq}}

Let $S=\C\setminus\{c_1,\dots,c_\mu\}$, $C_t=q^{-1}(t)$. The arguments below
use three geometric facts.
\begin{enumerate}[label=\textup{(G\arabic*)},leftmargin=2.8em]
\item $q$ is a locally trivial fibration over $S$; every fibre of $q$ is
 reduced and irreducible; for $t\in S$, $H_1(C_t;\C)$ has dimension $\mu$ and
 its intersection form is nondegenerate.
\item The local monodromy $T_i$ around $c_i$ is the Picard--Lefschetz
 transvection $\delta\mapsto\delta\pm\langle\delta,\nu_i\rangle\nu_i$, and the vanishing
 period $\int_{\nu_i(t)}dx\wedge dy/dq$ is holomorphic at $c_i$ with a nonzero
 value there; in particular $\nu_i\ne0$ and $T_i\ne1$.
\item $\pi_1(S)$ acts irreducibly on $H_1(C_t;\C)$.
\end{enumerate}

\begin{proposition}[Geometry of the class $(W)$]\label{prop:geometry}
For $q$ of class $(W)$ satisfying \textup{(M)}, \textup{(G1)--(G3)} hold.
\end{proposition}
\begin{proof}
\emph{Fibres.} The top weighted form of $q-t$ is $\alpha x^a+\beta y^b$, which is
irreducible in $\C[x,y]$ because $\gcd(a,b)=1$. The top form of a product is
the product of the top forms, so $q-t$ is irreducible; in particular every
fibre is reduced and irreducible.

\emph{Tameness.} Put $R=\max(|x|^{1/b},|y|^{1/a})$, $X=x/R^b$, $Y=y/R^a$. Then
$R^{-(ab-b)}q_x=a\alpha X^{a-1}+o(1)$ and $R^{-(ab-a)}q_y=b\beta Y^{b-1}+o(1)$
uniformly on $\max(|X|,|Y|)=1$ as $R\to\infty$, and one of the two leading terms
has modulus bounded below there. Hence $\|\nabla q\|\to\infty$ at infinity, so
$q$ is tame in the sense of Broughton \cite[Definition~3.1]{Br}: $\|\partial q\|$
is bounded away from $0$ outside a compact neighbourhood of the critical
points. For tame polynomials, \cite[p.~230]{Br} shows that $q$ is locally
trivial over any region of the target containing no critical value, and
\cite[Theorem~1.2 and p.~231]{Br} gives $H_1(C_t;\Z)\cong\Z^{\mu'}$ for a regular
value $t$, where $\mu'$ is the sum of the Milnor numbers at the critical points;
by Lemma~\ref{lem:normal}(c), $\mu'=\dim\C[x,y]/(q_x,q_y)=\mu$.

\emph{Local model and stability.} Write
\[
 f_\epsilon(X,Y)=\alpha X^a+\beta Y^b+
 \sum_{bi+aj<ab}q_{ij}\epsilon^{ab-bi-aj}X^iY^j.
\]
Every exponent of $\epsilon$ is a positive integer, so this is a holomorphic
deformation of $f_0=\alpha X^a+\beta Y^b$. For $\epsilon\ne0$ it equals
$\epsilon^{ab}q(\epsilon^{-b}X,\epsilon^{-a}Y)$; its critical points and values
are $(\epsilon^b x_i,\epsilon^a y_i)$ and $\epsilon^{ab}c_i$.
Choose a Milnor ball $\overline B_r$ for $f_0$, and choose $\eta>0$ so that its
levels $|t|\le2\eta$ are transverse to $\partial B_r$. Compactness and
$C^1$ convergence preserve this transversality for $f_\epsilon$ with
$|t|\le\eta$ when $|\epsilon|$ is sufficiently small. Shrink the parameter disk
so that all critical points lie in $B_{r/2}$ and all critical values have
modulus less than $\eta/2$.

At the fixed value $t_* =\eta$, the family
\[
 \{(\epsilon,X,Y):f_\epsilon(X,Y)=t_*,\ (X,Y)\in\overline B_r\}
 \longrightarrow\{\epsilon:|\epsilon|<\epsilon_0\}
\]
is a proper submersion, including on its boundary. The boundary form of
Ehresmann's theorem therefore identifies its fibres with the Milnor fibre of
$f_0$. This is also the perturbation comparison in \cite[\S2.1,
pp.~29--31]{AGV}. For a fixed small nonzero $\epsilon$, the local proper
submersion over the disk minus the critical values transports this
identification to every regular level there. Put
$S_t=f_\epsilon^{-1}(t)\cap\overline B_r$ and $C_t=f_\epsilon^{-1}(t)$.
By \cite[Theorem~2.1, p.~31]{AGV} a distinguished set of vanishing cycles
$\Delta_1,\ldots,\Delta_\mu$ is a basis of $H_1(S_t;\Z)$.
For a plane-curve singularity the kernel of the intersection form has
dimension $r_0-1$, where $r_0$ is the number of branches
\cite[\S3.6, p.~100]{AGV}. The binomial germ $f_0$ is irreducible, so this
form is nondegenerate; compare also \cite[p.~411, Example~2]{AGV}.

\emph{Compatible local and global transport.} Tameness is preserved by the
fixed nonzero rescaling. Over a relatively compact target neighborhood $U$
whose closure avoids the critical values, $\|\nabla f_\epsilon\|$ is bounded
below: use tameness outside a compact set and compactness inside it.
A bounded horizontal lift of a constant target tangent vector $v\in\C$ is
\[
 V_v=v\,\frac{\overline{\nabla f_\epsilon}}{\|\nabla f_\epsilon\|^2},
 \qquad df_\epsilon(V_v)=v.
\]
Boundary transversality supplies a smooth horizontal lift tangent to
$\partial B_r$. On a compact collar patch it to $V_v$ by a partition of unity.
The patched lift remains horizontal and bounded and is tangent to the sphere.
Its flows exist for the required finite times, preserve the sphere and its two
sides, and yield compatible local trivializations of the pair $(C_t,S_t)$.
This uses both tameness and transversality, following the bounded-lift
construction in \cite[pp.~229--230]{Br}.

Inclusion $j:S_t\hookrightarrow C_t$ preserves intersection numbers. Hence
$j_*v=0$ implies $\langle v,u\rangle=\langle j_*v,j_*u\rangle=0$ for every $u$,
and local nondegeneracy gives $v=0$. The global and local homology groups both
have dimension $\mu$, so $j_*$ is an isomorphism. The compatible transport
above makes it an isomorphism of local systems. The punctured disk contains
all critical values and its inclusion in their complement in $\C$ induces an
isomorphism on fundamental groups. Thus the global intersection form is
nondegenerate and its monodromy is that of this morsification. Rescaling back
to $q$ proves the remaining assertions of (G1).

\emph{(G2).} The Picard--Lefschetz formula is \cite[\S1.3, p.~26]{AGV}. At a
nondegenerate critical point choose holomorphic coordinates with
$q=c_i+uv$, $dx\wedge dy=J\,du\wedge dv$. The vanishing cycle is
$\{uv=t-c_i,\ |u|=r\}$ and $dx\wedge dy/dq=-J\,du/u$ on $C_t$, so the vanishing
period is $\mp2\pi i\sum_\kappa J_{\kappa\kappa}(t-c_i)^\kappa$, holomorphic with
value $\mp2\pi iJ(0)\ne0$ at $c_i$. Hence $\nu_i\ne0$, and $T_i\ne1$ by
nondegeneracy of the form.

\emph{(G3).} By the previous steps $H_1(C_t)$ has the basis
$\Delta_1,\dots,\Delta_\mu$, the global monodromy group is generated by the
Picard--Lefschetz transformations $T_{\Delta_i}$, and the form is nondegenerate.
By the theorem of Gabrielov and Lazzeri \cite[Theorem~3.3]{AGV}, the diagram of
a singularity relative to any distinguished basis is connected: the graph on
$\{1,\dots,\mu\}$ with an edge when $\langle\Delta_i,\Delta_j\rangle\ne0$ is
connected. (The diagram is defined after stabilization, which changes
off-diagonal intersection numbers only by signs; self-intersections may change.) Let $W\ne0$ be invariant. By
nondegeneracy there are $w\in W$ and $i$ with $\langle w,\Delta_i\rangle\ne0$;
then $(T_{\Delta_i}-1)w=\pm\langle w,\Delta_i\rangle\Delta_i$ puts $\Delta_i$ in
$W$, and $(T_{\Delta_j}-1)\Delta_i=\pm\langle\Delta_i,\Delta_j\rangle\Delta_j$
propagates membership along the connected graph. So $W=H_1$.
\end{proof}

For a horizontal family of cycles $\delta(t)\subset C_t$ over a simply connected
$U\subset S$ and $\omega=\sum_m z^m\omega_m\in\Omega^2_{\C[x,y]}[z]$ put
\[
 \pi_{\omega_m}(t)=\int_{\delta(t)}\frac{\omega_m}{dq},\qquad
 \Pi_\delta(\omega)(t)=\sum_m(-\partial_t)^m\pi_{\omega_m}(t).
\]
\begin{lemma}[Weyl intertwining]\label{lem:Pi}
$\Pi_\delta(z\omega)=-\partial_t\Pi_\delta(\omega)$,
$\Pi_\delta(\partial_z\omega+q\omega)=t\,\Pi_\delta(\omega)$, and
$\Pi_\delta(d_z\beta)=0$ for $\beta\in\Omega^1[z]$.
\end{lemma}
\begin{proof}
The first is the definition. Since $q=t$ on $C_t$, $\pi_{q\omega_m}=t\pi_{\omega_m}$,
and $(-\partial_t)^mt=t(-\partial_t)^m-m(-\partial_t)^{m-1}$ gives the second.
For a polynomial one-form $\beta_0$, $\pi_{dq\wedge\beta_0}=\int_{\delta(t)}\beta_0$
and $\pi_{d\beta_0}=\partial_t\int_{\delta(t)}\beta_0$ (Gelfand--Leray); the two
contributions to $\Pi_\delta((d+z\,dq\wedge)\beta)$ cancel termwise.
\end{proof}

\begin{lemma}[Period injectivity]\label{lem:injective}
Assume \textup{(G1)}, \textup{(G3)} and $\mu\ge2$. If $\Pi_\delta(\omega)\equiv0$
for every $\delta$, then $[\omega]=0$ in $\Hq\otimes\C$.
\end{lemma}
\begin{proof}
If $d\theta=\omega'$ then $d_z\theta=\omega'+z\,dq\wedge\theta$, so
$[\omega']=-z[dq\wedge\theta]$. Write $\omega=\sum_{m=0}^Nz^m\omega_m$ and apply this to the lowest coefficient:
subtracting $d_z\theta_0$ removes the $z^0$ term and changes only the
coefficient of $z^1$. After $N$ steps $\omega\equiv z^N\omega^\sharp$ modulo
$d_z\Omega^1[z]$ with $\omega^\sharp$ independent of $z$, and
$\Pi_\delta(\omega)=(-\partial_t)^N\pi_{\omega^\sharp}$ by Lemma~\ref{lem:Pi}. Thus
every branch of every period of $\alpha^\sharp=\omega^\sharp/dq$ is a polynomial of
degree $<N$ (zero when $N=0$), hence single valued: $\int_{(g-1)\delta}\alpha^\sharp=0$ for all
$g\in\pi_1(S)$. The span of the $(g-1)\delta$ is an invariant subspace, nonzero
because an irreducible representation of dimension $\ge2$ is nontrivial; by
(G3) it is all of $H_1$. So all periods of $\alpha^\sharp$ vanish. Choose
$\theta^\sharp$ with $d\theta^\sharp=\omega^\sharp$. Its periods have derivative
zero, hence are constant, hence single valued, hence (same argument) zero. So
$\theta^\sharp$ is exact along the generic fibres. The critical locus is finite
and all fibres are connected and nonempty; these give Bonnet's primitive
and quasi-fibered hypotheses for a map to the line. Thus
\cite[Definitions~1.3--1.4 and Theorem~1.5]{Bonnet} gives
$\theta^\sharp=dR+a\,dq$. Then
$\omega^\sharp=da\wedge dq=d_z(a\,dq)$ and $[\omega]=\pm z^N[\omega^\sharp]=0$.
\end{proof}

\begin{proposition}[Irreducibility]\label{prop:irred}
Assume \textup{(G1)--(G3)} and $\mu\ge2$. Every nonzero element of
$\Hq\otimes_{F_0}F$ has minimal differential operator of order $\ge\mu$ over
$F$. With Lemma~\ref{lem:normal}, $\Hq\otimes F$ is an irreducible differential
module of rank $\mu$.
\end{proposition}
\begin{proof}
Let $\eta\ne0$. Multiplying $\eta$ by an element of $F^\times$ changes neither
the differential submodule it generates nor the minimal order of an
annihilator (Leibniz), so we may assume $\eta=[\omega]$ with
$\omega\in\Omega^2_\C[z]$. Let
$P=\sum_{l=0}^rp_l(z)\nabla_z^l$, $p_l\in\C[z]$ not all zero, satisfy
$P\eta=0$. Multiplying $P$ on the left by a polynomial in $z$, which does not change $r$,
we may assume $P(z,\nabla_z)\omega=d_z\beta$ with $\beta\in\Omega^1[z]$. By
Lemma~\ref{lem:Pi}, for every $\delta$,
\[
 \widehat P\,\Pi_\delta(\omega)=0,\qquad
 \widehat P=\sum_lp_l(-\partial_t)\,t^l=\sum_{n=0}^Na_n(t)\partial_t^n,
 \qquad a_N(t)=(-1)^N\sum_l p_{l,N}\,t^l,
\]
where $N=\max\deg p_l$ and $p_{l,N}$ is the coefficient of $z^N$ in $p_l$. So
$a_N\ne0$ and $\deg a_N\le r$.
The map $\delta\mapsto\Pi_\delta(\omega)$ from $H_1(C_{t_0};\C)$ to germs at $t_0$
is linear, compatible with analytic continuation, and nonzero by
Lemma~\ref{lem:injective}; by (G3) it is injective. If $a_N(c_i)\ne0$, every
solution germ of $\widehat P$ near $c_i$ is holomorphic at $c_i$, so
continuation around $c_i$ is trivial on the image, contradicting $T_i\ne1$.
Hence $a_N$ vanishes at $c_1,\dots,c_\mu$ and $r\ge\deg a_N\ge\mu$. A proper
nonzero submodule would contain an element with minimal operator of order
$<\mu$.
\end{proof}

\section{The exponential torus and \texorpdfstring{$\SL_\mu$}{SL}}

Consider $Y'=AY$ with $A$ as in Lemma~\ref{lem:normal} and let $G\subset\GL_\mu$
be its differential Galois group over $F$, acting on the solution space $V$.

\begin{lemma}[Formal diagonalization]\label{lem:formal}
Let $A=\sum_{j\ge0}A_jw^j$, $w=1/z$, with $A_0$ having distinct eigenvalues
$c_1,\dots,c_n$. Then $Y'=AY$ has a formal fundamental matrix
$\widehat Y=\widehat H(w)\,z^{\Lambda}e^{Cz}$ with $C=\operatorname{diag}(c_i)$,
$\Lambda$ a constant diagonal matrix and $\widehat H\in\GL_n(\C[[w]])$.
\end{lemma}
\begin{proof}
After a constant change of basis $A_0=C$. Seek $\widehat H=\sum_{k\ge0}H_kw^k$,
$H_0=I$, with $\widehat H'=A\widehat H-\widehat H(C+\Lambda w)$; then
$\widehat Hz^\Lambda e^{Cz}$ is a solution. Since $(w^k)'=-kw^{k+1}$, the
coefficient of $w^k$ gives, for $k\ge1$,
\[
 -(k-1)H_{k-1}=[C,H_k]+\sum_{j=1}^kA_jH_{k-j}-H_{k-1}\Lambda .
\]
For $k=1$ put $\Lambda=\operatorname{diag}(A_1)$; the off-diagonal entries
$(c_i-c_l)(H_1)_{il}$ of $[C,H_1]$ determine the off-diagonal part of $H_1$.
For $k\ge2$ the diagonal part of the equation reads
$-(k-1)\operatorname{diag}H_{k-1}=\operatorname{diag}\bigl(A_1^{\rm off}H_{k-1}^{\rm off}
+\sum_{j=2}^kA_jH_{k-j}\bigr)$, which determines $\operatorname{diag}H_{k-1}$ from
quantities already known, and the off-diagonal part determines
$H_k^{\rm off}$. Induction on $k$ constructs $\widehat H$.
\end{proof}

\begin{proposition}\label{prop:SL}
Assume \textup{(M)}, \textup{(Q)}, \textup{(G1)--(G3)}, $\mu\ge2$. Then
$G\supseteq\SL_\mu$.
\end{proposition}
\begin{proof}
\emph{Formal solution and a diagonal subgroup.} Let $\mathfrak K=\C((w))$ with
$'=d/dz=-w^2d/dw$. By Lemma~\ref{lem:normal}(c) and Lemma~\ref{lem:formal},
$\widehat H'=A\widehat H-\widehat H(C+\Lambda w)$ with
$C=\operatorname{diag}(c_i)$, $\Lambda=\operatorname{diag}(\lambda_i)$,
$\widehat H\in\GL_\mu(\C[[w]])$. Let $K_0$ be a Picard--Vessiot extension of
$\mathfrak K$ for the diagonal system $Y'=(C+\Lambda w)Y$; it is generated by
nonzero $u_i$ with $u_i'=(c_i+\lambda_iw)u_i$. Then
$\widehat Y=\widehat H\operatorname{diag}(u_i)$ is a fundamental matrix of
$Y'=AY$ in $K_0$, and since $K_0$ has constants $\C$,
$\widehat L=F(\widehat Y_{ij})$ is a Picard--Vessiot extension of $F$; $G$ is its
group, written in the basis of columns of $\widehat Y$. Every
$\sigma\in\Gamma=\Gal(K_0/\mathfrak K)$ multiplies each $u_i$ by a constant, so
it restricts to an element of $G$ which is diagonal. Let $D_\Gamma\subset G$ be
the resulting closed subgroup of the diagonal torus. A character
$x^n=\prod x_i^{n_i}$ is trivial on $D_\Gamma$ iff $\prod u_i^{n_i}$ is
$\Gamma$-fixed, i.e.\ lies in $\mathfrak K$. An element $f\in\mathfrak K^\times$
has $f'/f\in w\C[[w]]$, while
$(\prod u_i^{n_i})'/\prod u_i^{n_i}=\sum n_ic_i+(\sum n_i\lambda_i)w$; so such
$n$ satisfy $\sum n_ic_i=0$. A closed subgroup of a torus is the common kernel
of the characters trivial on it, hence
\[
 D_\Gamma\supseteq T:=\{d:\ d^n=1\ \text{whenever}\ \textstyle\sum n_ic_i=0\}.
\]
$T$ is a torus with character group $\Z^\mu/\{n:\sum n_ic_i=0\}\cong\sum\Z c_i$,
acting on the $i$th basis vector with weight $c_i$; being connected,
$T\subset G^0$.

\emph{Root decomposition.} By (Q) the weights $c_i-c_j$ $(i\ne j)$ of $T$ on
$\mathfrak{gl}_\mu$ are pairwise distinct and nonzero, so
$\mathfrak g=\Lie G^0=(\mathfrak g\cap\mathfrak d)\oplus\bigoplus_{(i,j)\in R}\C E_{ij}$,
$\mathfrak d$ the diagonal matrices. $G$ is reductive because $V$ is a faithful
irreducible representation (Proposition~\ref{prop:irred}), so $G^0$ is
reductive; $\mathfrak g\cap\mathfrak d$ is the centralizer of $\Lie T$, a Cartan
subalgebra, and therefore $R=-R$. Also $R$ is closed under
$(i,j),(j,l)\mapsto(i,l)$. So $R$ defines an equivalence relation with classes
$I_1,\dots,I_s$; each $V_{I_a}=\Span\{e_i:i\in I_a\}$ is an irreducible
$G^0$-module and they are pairwise nonisomorphic, so the $G^0$-submodules of
$V$ are exactly the sums of the $V_{I_a}$.

\emph{One class.} Since $G^0$ is normal, $G$ permutes the $V_{I_a}$. The fixed
field $E_0$ of $G^0$ in $\widehat L$ is a finite Galois extension of $F$ with
group $G/G^0$. Under the isomorphism of $\widehat L$ with the Picard--Vessiot
field generated by an analytic fundamental matrix, its elements are algebraic
functions meromorphic on the universal cover of $\C^\times$ (the system is
singular only at $0,\infty$), so $E_0=F(\zeta)$ with $\zeta^d=z$. The image of
$D_\Gamma$ in $\Gal(E_0/F)$ has fixed field $E_0\cap\mathfrak K$. This is
$F(\zeta^{d/d'})$ for some $d'\mid d$, and
$(\zeta^{d/d'})^{d'}=w^{-1}$ forces $d'=1$ because $w$-orders in $\mathfrak K$
are integers. So $D_\Gamma$ maps onto $G/G^0$, i.e.\ $G=G^0D_\Gamma$. Diagonal
matrices preserve each $V_{I_a}$, so each $V_{I_a}$ is $G$-stable and
irreducibility gives $s=1$. Then every $E_{ij}$ lies in $\mathfrak g$ and
$\SL_\mu\subset G^0$.
\end{proof}

\section{Face disjointness}

\paragraph{Stable lattices at infinity.}
Put $R_\infty=\C[[w]]$, $K_\infty=\C((w))$ and
$\delta=d/dz=-w^2d/dw$. A lattice in a differential module over $K_\infty$
is a finite free $R_\infty$-submodule spanning it. A lattice $L$ is stable if
$\partial L\subset L$. Since $\delta(R_\infty)\subset wR_\infty$, its
operator induces a $\C$-linear operator on $L/wL$, called the leading
operator of this stable lattice. In the coefficient-module convention its
matrix for a system $Y'=E(w)Y$ is $E(0)^{\mathsf T}$, which has the same
spectrum as $E(0)$.

\begin{lemma}[Stable lattices and subquotients]\label{lem:lattice}
Let $L\subset M$ be a stable lattice and $N\subset M$ a differential
submodule. Then $L_N=L\cap N$ and $L_{M/N}=\operatorname{im}(L\to M/N)$
are stable lattices, and
\[
 0\longrightarrow L_N/wL_N\longrightarrow L/wL
 \longrightarrow L_{M/N}/wL_{M/N}\longrightarrow0
\]
is exact and compatible with the leading operators. Their characteristic
polynomials therefore multiply. In particular a direct sum of successive
quotients of a differential-module filtration has a stable lattice with the
same leading characteristic polynomial as $L$.
\end{lemma}
\begin{proof}
The intersection is a saturated, finitely generated submodule of $L$:
if $wv\in L\cap N$ and $v\in L$, then $v\in N$ because $N$ is a
$K_\infty$-vector space. The intersection and the quotient are free over the
discrete valuation ring $R_\infty$ and span the desired spaces. Stability
is inherited. Since the quotient is free, reduction modulo $w$ preserves
exactness. Also $\delta(w)=-w^2$, so each leading operator is well-defined
and the sequence intertwines them. Apply this successively to the filtration.
\end{proof}

\begin{lemma}[Detection of a rank-one leading coefficient]\label{lem:rankone}
If a module $M$ with stable lattice $L$ contains a nonzero vector $v$ with
$\partial v=(d+\rho w)v$, $d,\rho\in\C$, then $d$ is an eigenvalue of the
leading operator on $L/wL$.
\end{lemma}
\begin{proof}
The saturated lattice $L\cap K_\infty v$ has a generator $fv$, where
$f=w^m u$ with $u\in R_\infty^\times$. Its coefficient is
\[
 d+\rho w+\frac{\delta f}{f},\qquad
 \frac{\delta f}{f}=-mw-w^2\frac{u_w}{u}\in wR_\infty.
\]
Its leading eigenvalue is therefore $d$. Its reduction injects into $L/wL$
by Lemma~\ref{lem:lattice}.
\end{proof}

\begin{corollary}[A leading-spectrum obstruction]\label{cor:face-lattice}
Suppose $M$ has a stable lattice with leading spectrum $\Sigma$. Let $P$
admit a basis over $K_\infty$ with
$\partial e_i=(c_i+\lambda_iw)e_i$. If
\[
 \End^0(P)\hookrightarrow\End(M^{\rm ss})
\]
is a differential-module embedding, then $c_i-c_j\in\Sigma-\Sigma$ for
all $i\ne j$.
\end{corollary}
\begin{proof}
By Lemma~\ref{lem:lattice}, $M^{\rm ss}$ has a stable lattice $L^{\rm ss}$
with the same leading characteristic polynomial. The lattice
$\End_{R_\infty}(L^{\rm ss})$ is stable; its leading operator is the
commutator with that of $L^{\rm ss}$, with eigenvalues $s-s'$ for
$s,s'\in\Sigma$. The vector $e_i\otimes e_j^*$ is trace-free for $i\ne j$
and has differential coefficient
$c_i-c_j+(\lambda_i-\lambda_j)w$. Its image in the endomorphism module is
nonzero, so Lemma~\ref{lem:rankone} gives the result.

If $M^{\rm ss}$ was taken over $F$ first, extend a composition series from
$F$ to $K_\infty$ and apply the lattice lemma to this filtration. Its
quotients need not remain simple. No assertion that semisimplification
commutes with this extension is used.
\end{proof}

For a face polynomial $p$ of degree $d\ge2$, the one-variable reduction gives
a matrix in $M_{d-1}(k[z^{-1}])$ whose leading matrix is multiplication by
$p$ on $k[s]/(p')$. Thus its stable lattice has as leading spectrum the
critical values of $p$, also when those values repeat. For the interior
module, Lemma~\ref{lem:formal} supplies the diagonal basis required in
Corollary~\ref{cor:face-lattice}. These are all the local facts needed below;
no general formal-normal-form theorem is required.

Let $\Bcal$ be the closed homogeneous system satisfied by the boundary vector
(reduced face moments, endpoint exponentials, $1$), $\Tcal_\partial$ its tensor
category and $G_\partial$ its Galois group.

\begin{proposition}\label{prop:faces}
Suppose $G\supseteq\SL_\mu$, $\mu\ge2$, and that the interior module
has the formal diagonalization of Lemma~\ref{lem:formal}. If $\End^0(A)$
belongs to $\Tcal_\partial$, then there is a face $e$ with $d_e\ge2$ such that
every $c_i-c_j$ is a difference of two critical values of $p_e$. In particular
\textup{(F)} implies $\End^0(A)\notin\Tcal_\partial$.
\end{proposition}
\begin{proof}
The Galois group of $\End^0(A)$ is $\PGL_\mu$, so there is a surjection
$\rho:G_\partial\to\PGL_\mu$ through which $G_\partial$ acts on
$\mathfrak{sl}_\mu$ by $\operatorname{Ad}\circ\rho$. The unipotent radical dies,
so $\rho$ factors through $G'=\Gal(\Bcal^{\rm ss})\subset\prod_e\GL(W_e)\times(\text{torus})$,
$W_e$ the solution space of $H_{p_e}^{\rm ss}$. Write
$\Lie G'=\mathfrak z\oplus\mathfrak s_1\oplus\dots\oplus\mathfrak s_r$. Since
$\mathfrak{sl}_\mu$ is simple, $\ker d\rho$ contains $\mathfrak z$ and all simple
ideals but one, $\mathfrak s$; as $\ker d\rho$ is $\operatorname{Ad}(G')$-stable,
so is $\mathfrak s$, and $d\rho:\mathfrak s\to\mathfrak{sl}_\mu$ is a
$G'$-equivariant isomorphism. The projections
$\mathfrak s\to\mathfrak{gl}(W_e)$ are $G'$-equivariant and each is injective or
zero; they are not all zero because the remaining factor is abelian. For an
$e$ with injective projection, the $G_\partial$-module $\mathfrak{sl}_\mu$ embeds
in $\End(W_e)$. Using the covariant horizontal-fibre equivalence and the trace dualities
noted above, this gives a differential-module embedding
\[
 \End^0(A)\hookrightarrow\End(H_{p_e}^{\rm ss}).
\]
For the face module the leading spectrum is the set of critical values of
$p_e$. Corollary~\ref{cor:face-lattice}, after extension to $K_\infty$,
therefore makes every $c_i-c_j$ a difference of two such critical values.
This contradicts (F).
\end{proof}

\section{The saddle lemma}

Let $\Ncal=\Span_F\{1,\ \text{reduced face moments},\ \text{endpoint exponentials}\}$.

\begin{lemma}\label{lem:saddle}
Assume \textup{(G1)--(G3)}, $\mu\ge2$, and \textup{(S)} or \textup{(S$^*$)}. Then
$I_D(dx\wedge dy)\notin\Ncal$.
\end{lemma}
\begin{proof}
Suppose $P(z)I_D(dx\wedge dy)=\sum_jR_j(z)g_j$ with $0\ne P$, $R_j\in\C[z]$ and
$g_j$ among the generators. By Lemma~\ref{lem:phase-density}(b), the inverse Laplace transforms of
the face generators have algebraic densities on intervals free of the
finitely many face-critical, vertex and zero values; endpoint
exponentials and $1$ give point masses.  The compactly supported
distributions are determined by their Laplace transforms: restrict an
identically zero entire transform to the imaginary axis and apply Fourier
uniqueness. Hence on such an interval the analytic density
$\rho=\rho_{D,1}$ of $I_D(dx\wedge dy)$ satisfies:
$g=P(-\partial_t)\rho$ is algebraic over $\C(t)$. Under (S) this applies to
$(c,c+\varepsilon)$ for small $\varepsilon>0$.

\emph{Local form under (S).} By the real-analytic Morse lemma there are local
coordinates with $q=c+uv$, $dx\wedge dy=J\,du\wedge dv$,
$J(0)^2=|\det\operatorname{Hess}q(P_0)|^{-1}$. Take a box
$\{|u|,|v|\le\delta\}$ on which the multiplicity of $D$ is $w$. After shrinking $\delta$ the edges of the box are transverse to the levels near
$c$ and its corners have levels $c\pm\delta^2\ne c$. Outside the open
box the level curves near level $c$ meet neither a critical point (distinct
critical values) nor a tangency or vertex of $\partial D$ (by (S)), so that
contribution to $\rho$ is analytic at $c$. Inside, with $s=t-c$,
\[
 \rho_{\rm box}(s)=\pm w\int_{|s|/\delta\le|u|\le\delta}J(u,s/u)\frac{du}{|u|}
 =\pm w\bigl(-2j(s)\log s+\text{analytic}\bigr)\quad(0<s<\delta^2),
\]
where $J=\sum J_{\kappa\lambda}u^\kappa v^\lambda$ and $j(s)=\sum_\kappa J_{\kappa\kappa}s^\kappa
=\frac1{2\pi i}\oint_{|u|=r}J(u,s/u)\frac{du}u$; the terms with
$\kappa\ne\lambda$ integrate to convergent power series in $s$. Since
$dx\wedge dy/dq=-J\,du/u$ on $C_t$, $j$ is $\mp\frac1{2\pi i}$ times the period
$h(t)=\int_{\nu(t)}dx\wedge dy/dq$ over the vanishing cycle, and $j(0)=J(0)\ne0$.
Thus $\rho$ continues around $c$ with variation a nonzero constant times $h$,
which is (S$^*$).

\emph{Conclusion from (S$^*$).} Write the variation as $\lambda h$, $\lambda\ne0$. The
vanishing period $h$ is fixed by the local monodromy, so $N$ turns around
$c_i$ change $\rho$ by $N\lambda h$. The germ $g$ continues with $\rho$ and
remains algebraic, so for some $N\ge1$ it returns to itself; therefore
$N\lambda P(-\partial_t)h=0$ near $c_i$ and $h$ extends to an entire function. So $h$ is
single valued on $S$, i.e.\ $\int_{(g-1)\nu_i}dx\wedge dy/dq=0$ for all
$g\in\pi_1(S)$. As in Lemma~\ref{lem:injective} the $(g-1)\nu_i$ span $H_1$
($\nu_i\ne0$ is not invariant, by (G3) and $\mu\ge2$), so the period over $\nu_i$
itself vanishes, contradicting (G2).
\end{proof}

\section{Proof of Theorem \ref{thm:main}}

(i) is Lemma~\ref{lem:normal}. For (ii): integrating the certificates of
$qm$ gives $M'=AM+b$ with $b_m=\int_{\partial D}e^{zq}(P_{qm}\,dy-Q_{qm}\,dx)\in\Ncal$.
The actual one-column criterion, Theorem~\ref{thm:criterion}, requires:
(1) $G\supseteq\SL_\mu$, which is Proposition~\ref{prop:SL};
(2) $\End^0(A)\notin\Tcal_\partial$, which is Proposition~\ref{prop:faces}
with (F) and Lemma~\ref{lem:formal};
(3) some $M_m\notin\Ncal$, which is Lemma~\ref{lem:saddle} since $1\in B$.
Hence the $M_m$ are algebraically independent over $F(\Ncal)=K_\partial$.

(iii) The vector (boundary coordinates, $M$) consists of $E$-functions
(Theorem~\ref{thm:boundary-system}) and satisfies a homogeneous linear system over $k(z)$
whose only finite pole is $z=0$: for $M$ by Lemma~\ref{lem:normal}, for the
faces by its one-variable case. Let $d$ be the transcendence degree of
$K_{\partial,0}$ over $k(z)$.  Since
$K_{\partial,0}\subset K_\partial=\C K_{\partial,0}$, the algebraic
independence in part (ii) implies algebraic independence over
$K_{\partial,0}$ itself.  Hence the combined boundary-and-moment field
has transcendence degree $d+\mu$ over $k(z)$. Beukers' form of Siegel--Shidlovskii
\cite[Theorem~1.1, p.~370]{Beukers}, applied first to the boundary vector
and then to the combined vector, gives the numerical degrees $d$ and
$d+\mu$ at every $\xi\in k^\times$. Their difference is $\mu$.
The one-variable normal forms have only powers of $z$ as denominators, so
these reduced boundary values generate the stated face-value field at
all such $\xi$.

(iv) Specialize (i) at $z=\xi$: the boundary term lies in $K_{\partial,\xi}$,
so algebraicity forces all $a_{r,m}(\xi)=0$ by (iii), which gives the exactness;
conversely exactness gives membership by Stokes. \qed

\section{The rank-six example}\label{sec:example}

Let $q=x^4+y^3+xy-y$, so $(a,b)=(4,3)$, $\mu=6$,
$B=\{1,x,x^2,y,xy,x^2y\}$, and $D$ the triangle of Example~\ref{ex:intro}.
The following are exact computations (script supplied).
\begin{itemize}[leftmargin=1.6em]
\item \emph{Connection.} $A=A_0+A_1/z$ with
 $\det(t-A_0)$ proportional to
 \[
 \chi(t)=1289945088\,t^6-573308928\,t^4+83980800\,t^3+83854656\,t^2+5477429\,t-4197429.
 \]
 The same polynomial is obtained independently as the resultant eliminating
 the critical points.
\item \emph{(M) and (Q).} $\chi$ is irreducible over $\Q$; modulo $7$ it
 factors with degrees $(1,5)$ and modulo $31$ with degrees $(1,2,3)$. A
 transitive group of degree $6$ containing a $5$-cycle is primitive, and with
 a transposition (the cube of the second Frobenius) it is $S_6$. Given a
 relation $\sum a_i c_i=0$ with rational $a_i$, applying a transposition
 and subtracting yields $(a_i-a_j)(c_i-c_j)=0$. Thus all $a_i$ are equal;
 a relation among the differences has coefficient sum zero and is trivial.
\item \emph{(S).} The critical points are $x=1-3y^2$ with
 $-108y^6+108y^4-36y^2+y+4=0$. Write $f(y)$ for this sextic. Sturm counts give exactly two real roots, one in
 $(-\tfrac{420}{1000},-\tfrac{419}{1000})$ and one in
 $(\tfrac{722}{1000},\tfrac{723}{1000})$. On the critical locus
 $\det\operatorname{Hess}q=72x^2y-1=-f'(y)$, and $f'$ has no root in either
 interval, with $f'>0$ on the first and $f'<0$ on the second. So the first root
 gives a saddle $P_0\approx(0.471644,-0.419665)$, $c=q(P_0)\approx0.197305$,
 and the second does not. The polynomials $x-\tfrac{y+1}2$ and
 $1-\tfrac{y+1}2-x$ with $x=1-3y^2$ have no root on the first interval and are
 positive there, so $P_0$ is interior to $D$.
\item \emph{(F) and the rest of (S)} hold by Corollary~\ref{cor:polygon} with
 $k_0=\Q$, since $[\Q(c):\Q]=6\ge4$.
\end{itemize}
For reproducibility, writing $\bar\chi_p$ for the monic reduction modulo $p$,
\begin{align*}
 \bar\chi_5&=t^6-t^4+2t^2-2t+2,\\
 \bar\chi_7&=(t+3)(t^5-3t^4+3t^2+3t+2),\\
 \bar\chi_{31}&=(t-14)(t^2+8t+10)(t^3+6t^2-14t+1).
\end{align*}
The displayed factors modulo $7,31$ are irreducible and squarefree.
Irreducibility modulo $5$ is certified by
\[
 \gcd(\bar\chi_5,t^{5^2}-t)=\gcd(\bar\chi_5,t^{5^3}-t)=1,
 \qquad t^{5^6}\equiv t\pmod{\bar\chi_5}.
\]
All these are exact polynomial computations in the accompanying certificate.
The supplementary exact checks give
$\operatorname{tr}(A_0^k(A_1+I))=0$ for $0\le k\le5$.
Since $A_0$ has simple spectrum, the Vandermonde matrix in its eigenvalues
shows that all six diagonal entries of $A_1$ in its eigenbasis are $-1$.
Thus the six powers in Lemma~\ref{lem:formal} are exactly $-1$.

The exact checks are run by \texttt{make test}. No numerical quadrature
or asymptotic fitting is used to certify the hypotheses.

\section{The elliptic triangle and the continuation alternative}\label{sec:elliptic-example}
Fix the counterclockwise triangle
\[
 D=\operatorname{conv}\{(-1/2,-1),(1/2,0),(-1/2,1)\}
 =\{-1/2\le x\le1/2,\ |y|\le1/2-x\},
\]
and
\[
 q=y^2-x^3+3x,\qquad
 M_j=\int_Dx^j e^{zq}\dd x\dd y\quad(j=0,1).
\]
A polynomial parametrization is
$\Phi(u,v)=(u-1/2,(1-u)(2v-1))$, with Jacobian $2(1-u)$.

For $\beta=P\,dy-Q\,dx$, one has
\[
 d_z\beta=(L_xP+L_yQ)\,dx\wedge dy,
\quad L_x=\partial_x+3z(1-x^2),\quad L_y=\partial_y+2zy.
\]
The operators act in separate variables. Their one-variable quotients have
bases $1,x$ and $1$, respectively; both operators are injective on
polynomials. Tensoring the two one-variable complexes proves
\begin{equation}\label{eq:ranktwo}
 \Hq=F_0[dx\wedge dy]\oplus F_0[x\,dx\wedge dy].
\end{equation}
In particular the rank is exactly two, not merely at most two.

The explicit reduction recurrences are
\begin{align}
 [x^i y^j]&=-\frac{j-1}{2z}[x^i y^{j-2}] &&(j\ge1),\label{eq:yrec}\\
 [x^{m+2}]&=[x^m]+\frac{m}{3z}[x^{m-1}] &&(m\ge0),\label{eq:xrec}
\end{align}
with the zero-coefficient terms omitted. Each use records a polynomial
primitive. Hence every $r\in k[x,y]$ has a unique remainder
$a_r(z)+b_r(z)x$, with $a_r,b_r\in k[z,z^{-1}]$, and an identity
\begin{equation}\label{eq:elliptic-normal}
 r=a_r+b_rx+L_xP_r+L_yQ_r,
 \qquad P_r,Q_r\in k[z,z^{-1},x,y].
\end{equation}

Let
\[
 C=\begin{pmatrix}0&2\\2&0\end{pmatrix},\qquad
 D_0=\begin{pmatrix}-5/6&0\\0&-7/6\end{pmatrix},\qquad A=C+D_0/z.
\]
The following are identities of polynomial forms, with Laurent coefficients
in $z$:
\begin{align}
 q\,dx\wedge dy
 &= -\frac5{6z}\,dx\wedge dy+2x\,dx\wedge dy+d_z\beta_0,\\
 qx\,dx\wedge dy
 &=2\,dx\wedge dy-\frac7{6z}x\,dx\wedge dy+d_z\beta_1,\label{eq:explicit-beta}
\end{align}
where
\[
 \beta_0=\frac{x}{3z}\,dy-\frac{y}{2z}\,dx,
 \qquad
 \beta_1=\frac{x^2-2}{3z}\,dy-\frac{xy}{2z}\,dx.
\]
Define actual face integrals
\[
 B_{x,j}=\int_{\partial D}x^je^{zq}\,dy,
 \qquad
 B_{y,j}=-\int_{\partial D}x^jye^{zq}\,dx.
\]
Then
\begin{equation}\label{eq:elliptic-system}
 \binom{M_0}{M_1}'
 =\begin{pmatrix}-5/(6z)&2\\2&-7/(6z)\end{pmatrix}\binom{M_0}{M_1}
 +\frac1z\binom{B_{y,0}/2+B_{x,1}/3}
 {B_{y,1}/2+B_{x,2}/3-2B_{x,0}/3}.
\end{equation}
This is obtained by integrating \eqref{eq:explicit-beta}, not by choosing
an isomorphic abstract connection and then normalizing a solution.

\subsection{A finite actual boundary field}
Put
\[
 p(x)=-x^3+x^2+2x+\tfrac14,
 \quad P_j(z)=\int_{-1/2}^{1/2}x^j e^{zp(x)}\dd x,
 \quad Q_0(z)=\int_{-1}^{1}e^{z(y^2-11/8)}\dd y,
\]
and $E_a=e^{-3z/8}$, $E_b=e^{11z/8}$. Parametrizing the three oriented
edges gives
\begin{equation}\label{eq:face-identities}
 B_{x,j}=2P_j-(-1/2)^jQ_0,
 \qquad B_{y,j}=P_j-2P_{j+1}.
\end{equation}
The one-variable reductions imply
\begin{equation}\label{eq:boundary-field}
 K_{\partial,0}=k(z)(P_0,P_1,Q_0,E_a,E_b).
\end{equation}
Indeed the two slanted faces have the same phase $p$; their arbitrary
polynomial amplitudes reduce to $P_0,P_1$ and endpoints. The vertical
quadratic amplitudes reduce to $Q_0$ and endpoints.

For completeness, a closed boundary system is
\begin{align}
 P_0'&=\left(\frac{17}{36}-\frac1{3z}\right)P_0+
       \frac{14}{9}P_1+\frac{E_b+5E_a}{18z},\\
 P_1'&=\left(\frac{28}{27}+\frac1{9z}\right)P_0+
       \left(\frac{163}{108}-\frac2{3z}\right)P_1+
       \frac{-53E_b+41E_a}{108z},\\
 Q_0'&=\left(-\frac{11}{8}-\frac1{2z}\right)Q_0+\frac{E_a}{z},
 \qquad E_a'=-\tfrac38 E_a,\quad E_b'=\tfrac{11}{8}E_b.
\end{align}
For example, $P_2=(2P_1+2P_0)/3-(E_b-E_a)/(3z)$, so
\eqref{eq:elliptic-system} closes with this boundary vector. The only
finite pole of the combined rational system is zero.

\subsection{The actual normalization at zero}
The face forcing in \eqref{eq:elliptic-system} is $b(z)/z$ with $b$ entire.
Writing $M=\sum v_mz^m$, $b=\sum b_mz^m$, its holomorphic solution obeys
\[
 (mI-D_0)v_m=Cv_{m-1}+b_m,\qquad v_{-1}=0.
\]
Every $mI-D_0$, $m\ge0$, is invertible. Thus there is at most one
holomorphic solution at zero with these actual forcing functions. The
integral supplies it, with
\[
 M_0(0)=1,\qquad M_1(0)=-\tfrac16.
\]
This supplies a second, independent check that no unspecified homogeneous
solution has been added to the actual column.

\subsection{The actual continuation}
On the regular phase fibre $y^2=x^3-3x+t$, the convention
$dq\wedge\alpha=dx\wedge dy$ gives $\alpha=-dx/(2y)$.

\begin{lemma}[Continuation for the elliptic triangle]\label{lem:elliptic-continuation}
The density germ of $M_0$ just to the right of $t=0$ satisfies $(S^*)$.
\end{lemma}
\begin{proof}
We verify a nonzero elliptic monodromy variation of the actual density germ.
Put $f(x)=-x^3+3x$, let $r_2(s)\in(-1,1)$ be its middle inverse branch
for $-2<s<2$, and let $c(t)$ be the inverse branch of $p$ beginning at
$c(0)\in(-1/5,0)$. Set $a(t)=1/2-c(t)$. Coarea on the triangle gives,
near zero,
\begin{equation}\label{eq:actual-density}
 \rho(t)=\int_{-a(t)}^{a(t)}
 \frac{dy}{3(1-r_2(t-y^2)^2)}.
\end{equation}
Here $p'(x)=-3x^2+2x+2$ is positive on $[-1/5,1]$.
Moreover $p(-1/5)<0<p(0)$ and $p(4/5)<2<p(1)$, so $c(t)$ is a
single real-analytic branch for $0\le t\le2$, with
$c(0)\in(-1/5,0)$ and $c(2)\in(4/5,1)$. For
$|y|\le|a(t)|$ one has
\[
 f(c(t))=t-a(t)^2\le t-y^2\le t<2,
 \qquad f(c(t))>-2.
\]
Thus the middle inverse branch is regular on this moving segment for
$0<t<2$. Writing the signed integral on the fixed interval as
\[
 a(t)\int_{-1}^1
 \frac{ds}{3(1-r_2(t-a(t)^2s^2)^2)}
\]
proves real analyticity and hence local holomorphic continuation along
that interval. At $t=11/8$ one has $c(t)=1/2$, so the endpoints meet,
but the displayed expression is analytic and zero there. This specifies
the continuation through that otherwise excluded boundary value. Near
$t=2$, $a(t)<0$ and $a(2)\ne0$, and both endpoints are unramified.

Write $\varepsilon=2-t$ and $r_2(t-y^2)=1-v$. Since
$f(1-v)=2-3v^2+v^3$, near the node the integrand is
\[
 \frac1{2\sqrt3\sqrt{\varepsilon+y^2}}+O(1).
\]
Choose a fixed $b>0$ smaller than $|a(2)|$. The outer parts
$|y|\ge b$ extend holomorphically to $t=2$; on $|y|\le b$ the remainder
in the displayed expansion is uniformly bounded for $\varepsilon>0$.
Since the integration orientation is reversed near $2$, it follows that
\begin{equation}\label{eq:audited-log-coefficient}
 \rho(t)=\frac{1}{2\sqrt3}\log(2-t)+O(1)
 \qquad(t\to2^-).
\end{equation}
The local geometry makes the corresponding variation precise. Put
$v=1-x$ and $u=v\sqrt{3-v}$. Then $q-2=y^2-u^2$ near the node, so a
nearby local fiber is an annulus. The path outside a small nodal
neighborhood and its two endpoints return under a sufficiently small
loop. Its variation lies in the one-dimensional cycle space of that
annulus. The coefficient in \eqref{eq:audited-log-coefficient} proves it
is nonzero. The local core cycle $\delta_+$ is fixed by this nodal
monodromy (equivalently by its Dehn twist). Consequently
\[
 (T_+-I)\rho=c_*h_+(t),\qquad c_*\ne0,
 \quad h_+(t)=\int_{\delta_+}\frac{dx}{2y},
 \quad T_+h_+=h_+.
\]
This is also the local annulus proof of the Picard--Lefschetz formula in
this case; the nonzero coefficient has just been checked on the actual
path by \eqref{eq:actual-density}.

The variation is a nonzero constant multiple of the required
Gelfand--Leray vanishing period; the common sign is immaterial.
\end{proof}

\begin{remark}
This is continuation of a regular density germ, not a claim that the
original compactly supported density is nonzero outside $q(D)=[-11/8,11/8]$.
\end{remark}

\subsection{Verification of the comparison hypotheses}
The phase has $(a,b)=(3,2)$ and its two critical values are $-2,2$.
Thus (M) holds and (Q) is automatic. For a slanted edge put
$\beta=7\sqrt7/27$. The substitution $x=1/3+(\sqrt7/3)s$ gives
\[
 p(x)=\frac{107}{108}+\beta(-s^3+3s).
\]
Its critical values are $107/108\pm2\beta$, so their nonzero differences
are $\pm4\beta$, whereas the interior differences are $\pm4$.
Since $\beta^2=343/729\ne1$, these do not agree. The vertical quadratic
face has only one critical value, so it also satisfies (F).
Lemma~\ref{lem:elliptic-continuation} supplies $(S^*)$, although the saddle
$(1,0)$ lies outside the triangle and (S) fails.

\begin{corollary}[Elliptic functional and numerical comparison]
\label{thm:intro-elliptic}\label{thm:elliptic-main}\label{thm:values}
For this triangle and phase, $M_0,M_1$ are algebraically independent over
$\C(z)(P_0,P_1,Q_0,E_a,E_b)$. At every nonzero algebraic $\xi$, their values
are algebraically independent over the corresponding fixed-face value
field. Every polynomial amplitude has the unique remainder
\eqref{eq:elliptic-normal}, and its value is algebraic over the face-value
field exactly when its specialized polynomial twisted class is zero.
\end{corollary}
\begin{proof}
Apply Theorem~\ref{thm:main}, using the verifications above. The formulas
\eqref{eq:explicit-beta}--\eqref{eq:boundary-field} retain the actual
normalization and the finite generating boundary vector. Thus this is
also a specialization of the general proof, not just an abstract
identification of two homogeneous equations.
\end{proof}
\section{Scope, dependencies, and verification}\label{sec:scope}
The normal-form theorem is algebraic and the polynomial boundary-system
construction allows arbitrary polynomial phases. Algebraic independence
requires the explicitly stated hypotheses of Theorem~\ref{thm:main},
verified for the two examples above. The result is not contingent on
an unproved genericity conjecture: (Q) is a hypothesis on the specified
critical values, certified in the examples. Its failure alone does not
imply that the homogeneous Galois group is smaller.

All boundary fields are attached to the finitely many fixed faces and
endpoints of the chain. Replacing them by the field of all compact
polynomial-line periods is not an asserted extension. Only one phase and
one actual column are considered. No several-chain or multiblock formal
completeness is inferred. Numerical independence is a statement after
specializing the parameter; it is not injectivity of evaluation on the
unspecialized algebra containing $z-\xi$.

\paragraph{External inputs.}
The geometric proof uses Broughton \cite[Definition~3.1, Theorem~1.2,
pp.~229--231]{Br} and the Picard--Lefschetz, distinguished-basis,
connectedness, and nondegeneracy statements in
\cite[\S1.3; \S2.1; Theorem~3.3; pp.~100, 411]{AGV}.
The needed perturbation and compatible-transport arguments are included
in Proposition~\ref{prop:geometry}. Polynomial relative exactness uses
\cite[Definitions~1.3--1.4, Theorem~1.5]{Bonnet}.
Ordinary Picard--Vessiot theory and its tensor equivalence use
\cite[Theorem~1.5.2]{Singer}, and the complete reducibility in the actual
criterion uses \cite[Proposition~15.15, Corollary~22.43]{MilneAG}.
The general finiteness assertion uses the explicitly identified
Proposition~9.1 of Sabbah \cite{Sabbah}.
Numerical specialization uses \cite[Theorem~1.1]{Beukers} separately for
the boundary and combined vectors. Neither Wasow's formal simplification
nor a general formal-normal-form theorem is needed: the relevant
simple-spectrum and lattice statements are proved in the text.

\paragraph{Computational verification.}
The repository contains the retained exact elliptic-amplitude certificate,
the rank-six polynomial/resultant/Frobenius/Sturm certificate, and exact
checks of the formal recursion and six formal powers. The rank-six
geometric inequalities are certified on whole rational intervals.
These tests verify finite calculations, not the general monodromy or
Galois proofs. Compilation, regression tests, and AI-assisted audits are
not independent specialist refereeing or formal verification.

\paragraph{AI assistance and responsibility.}
OpenAI ChatGPT and Anthropic Claude were used for proof exploration,
drafting, symbolic checks, and adversarial review. They are not authors.
The authors are responsible for the statements, proofs, citations, and
remaining errors.

\small
\begin{thebibliography}{99}

\bibitem{AGV} V.~I.~Arnold, S.~M.~Gusein-Zade, A.~N.~Varchenko,
 \emph{Singularities of Differentiable Maps}, vol.~II, Birkh\"auser, 1988.

\bibitem{Beukers} F.~Beukers,
 \emph{A refined version of the Siegel--Shidlovskii theorem}, Ann. of Math.
 \textbf{163} (2006), 369--379.
 \href{https://doi.org/10.4007/annals.2006.163.369}{doi:10.4007/annals.2006.163.369}.

\bibitem{Bonnet} P.~Bonnet,
 \emph{Relative exactness modulo a polynomial map and algebraic
 $(\mathbb C^p,+)$-actions}, Bull. Soc. Math. France \textbf{131} (2003),
 no.~3, 373--398. \href{https://doi.org/10.24033/bsmf.2447}{doi:10.24033/bsmf.2447}.

\bibitem{Br} S.~A.~Broughton, \emph{Milnor numbers and the topology of
 polynomial hypersurfaces}, Invent.\ Math.\ \textbf{92} (1988), 217--241.

\bibitem{CommelinHabeggerHuber}J.~Commelin, P.~Habegger, and A.~Huber,
\emph{Exponential periods and o-minimality},
preprint, arXiv:2007.08280, version~4, 2025.
\href{https://arxiv.org/abs/2007.08280}{arXiv:2007.08280}.

\bibitem{KontsevichZagier}M.~Kontsevich and D.~Zagier,
\emph{Periods},
in \emph{Mathematics Unlimited--2001 and Beyond},
Springer, Berlin, 2001, pp.~771--808.
\href{https://doi.org/10.1007/978-3-642-56478-9_39}{doi:10.1007/978-3-642-56478-9\_39}.

\bibitem{MilneAG}J.~S.~Milne,
\emph{Algebraic Groups: The Theory of Group Schemes of Finite Type
 over a Field}, Cambridge Studies in Advanced Mathematics, vol.~170,
Cambridge University Press, 2017, Proposition~15.15 and Corollary~22.43.
\href{https://www.jmilne.org/math/Books/iAG2017.pdf}{Author-hosted text}.

\bibitem{PaperVPrevious} C.~D.~Long and A.-A.~Le~Blanc,
 \emph{Relative Exponential Periods on Polynomial Flags:
 Polynomial Boundary Systems and an Elliptic Comparison Theorem},
 Paper~V, September 2026; repository snapshot
 \href{https://github.com/long-mathematics/exponential-integral-theorem/tree/d5ae8f9081c7bc34d1622316bd8463482f244395}{d5ae8f90},
 inspected 6 October 2026.

\bibitem{Sabbah}C. Sabbah, \emph{Hypergeometric periods for a tame polynomial},
Portugal. Math. \textbf{63} (2006), 173--226.
Preprint \href{https://arxiv.org/abs/math/9805077}{arXiv:math/9805077}.
We use Proposition~9.1 of the
\href{https://perso.pages.math.cnrs.fr/users/claude.sabbah/articles/sabbah_bbases.pdf}{author-hosted manuscript};
no result under its subsequent standing tameness assumption is used
for general finiteness.

\bibitem{Singer} M.~F.~Singer,
 \emph{Introduction to the Galois theory of linear differential equations},
 \href{https://arxiv.org/abs/0712.4124}{arXiv:0712.4124v2}, 2008.
 We use ordinary Picard--Vessiot theory and Theorem~1.5.2, not the
 general formal-normal-form theorem of \S1.4.

\end{thebibliography}
\end{document}
